\section{Prime bounds and the finite range}
\label{sec:arithmetic}

Take \(a,d\ge1\), \(N\ge25\), \(d<3a\), \(6a<(5N-11)d\), and
\(s=d^2/(a+d)<\sigma_N\), as in Lemma~\ref{lem:support-reduction}.
In this section \(m\) runs through \(N-1,N\), \(b=a+d\), \(M=bN\),
and \(F\) is a nonconstant monic rational common divisor of
\(R_m(x)=x^mP_m(1/x)\). By Gauss's lemma \(F\in\Z[x]\), and \(F(0)\ne0\).
Write \(h=\deg F\) and \(H=dh\). The reciprocal of even order is
positive on \(\R\), as is seen by integrating \((x+u^d)^m\) against
the gamma density of shape \(am+1\). Hence \(h\) is even.

In fact \(h<N-1\). To see this, put \(n=N-1\) and \(u=an+1\).
If \(P_n\mid P_{n+1}\), the quotient is \(1+ct\). Its linear
coefficient would have to be both
\[
 c_1=(n+1)(u+a)_d-n(u)_d,\qquad
 c_2=\frac{(u+dn)_b}{(u)_a}.
\]
Set \(A=(u)_b\), \(B=(u)_a(u)_d\), and \(f(z)=(u+z)_b\).
Then \((u)_a(c_2-c_1)=f(dn)-(n+1)A+nB\).
Convexity gives \(f(dn)\ge nf(d)-(n-1)A\), and
\(f(d)/A=(u+b)_d/(u)_d>A/B\). Thus
\(f(d)-2A+B>(A-B)^2/B\ge0\), proving \(c_2>c_1\).

A sufficiently large prime can occur in a common root valuation only
through a multiple close to the end of each factorial product. The
locations of those multiples must agree to within one coefficient step
after translation by \(b\). We first use this observation to restrict
primes in the two initial factorial blocks.

For a prime \(p\) put \(y_m(k)=v_p([t^k]P_m)\) and
\[
 U_m(j)=y_m(m)-y_m(m-j)
       =v_p\!\left(\prod_{i=0}^{dj-1}(bm-i)\right)-v_p\binom mj .
\]
A lower Newton slope \(\lambda\) of \(P_m\) is a valuation of roots of
\(R_m\). Its width is the number of roots at that valuation, counted
with multiplicity. Equivalently the extremes of the maximizing indices
of \(U_m(j)-\lambda j\) have this difference.

\begin{lemma}[Degree allocations at a prime]
\label{lem:newton-alloc}
For a common slope with reduced denominator \(v\), let \(w\) be the
smaller of its two widths. The common factor can use only a number of
roots from \(\{0,v,2v,\ldots,w\}\) at that slope. Consequently \(h\)
belongs to the sum of these sets over the common slopes, at every prime.
In particular each positive valuation of a root of \(F\) is at least
\(1/h\).
\end{lemma}
\begin{proof}
At each slope the roots of a common factor occur with at most the
smaller multiplicity. Its Newton polygon has integral vertices, so
each width in that polygon is divisible by the reduced denominator
of the slope. The last assertion follows since such a width is at
most \(h\).
\end{proof}

\begin{lemma}[Prime support of short blocks]
\label{lem:prime-support}
Set
\[
 C_N=\max\{a+2d-1,\sqrt M\},\qquad B_h=\min\{dN,\max(2H,C_N)\}.
\]
Every prime divisor of \(F(0)\) or of either block
\(\prod_{i=1}^H(am+i)\) is at most \(B_h\).
\end{lemma}
\begin{proof}
If \(p>dN\), then \(p^2>M\) because \(a<d(N-2)\), and all binomial
coefficients are \(p\)-adic units. Each full factorial interval
\((am,bm]\) has at most one multiple \(u_m\) of \(p\). If present,
it yields only one positive slope, equal to \(1/w_m\) with
\(w_m=1+\lfloor(bm-u_m)/d\rfloor\).
A common positive slope would require the same \(w_m\) in both orders.
It would give
\begin{equation}
 u_N-u_{N-1}\in[b-d+1,b+d-1]=[a+1,a+2d-1].
 \label{eq:translated-multiples}
\end{equation}
This interval is positive and below \(p\), so it contains no multiple
of \(p\).

Now let \(p>\max(2H,C_N)\) and suppose that \(F\) has a positive slope
\(\lambda\) at \(p\). By Lemma~\ref{lem:newton-alloc},
\(\lambda\ge1/h\). Since \(p^2>M\), the count of terminal multiples
gives \(U_m(j)<1+dj/p\) for \(j>0\). The maximizing value of
\(U_m(j)-\lambda j\) is at least zero, so every maximizing index is
less than \(2h\). Throughout this range \(j<p\) and \(dj<p\).
The binomial valuation cannot decrease because \(j!\) is a unit.
The factorial valuation can rise only once and by one. The weighted
ordinate therefore decreases strictly between jumps. A positive
slope must tie zero to the index of the single factorial jump, giving
slope \(1/w_m\) with \(w_m=1+\lfloor(bm-u_m)/d\rfloor\).
As before, a shared slope implies \eqref{eq:translated-multiples},
which is impossible since \(p>C_N\).

These two cases show that \(p>B_h\) cannot divide \(F(0)\).
If it divides one of the first \(H\) factorial terms, all the
coefficients at \(k\ge h\) in \(P_m\) vanish modulo \(p\).
Then \(R_m\bmod p\) is a power of \(x\) times a polynomial of degree
at most \(h-1\). Its monic divisor \(F\bmod p\) has degree \(h\)
and nonzero constant term, a contradiction.
\end{proof}

\subsection{The degree bound}

\begin{lemma}[Degree bound from the first factorial block]
\label{lem:degree-cap}
If \(aN+1>\max\{H+13,279H/262\}\), then \(2H<C_N\).
The starting-point inequality holds whenever \(a\ge2d\), or
whenever \(N\ge275\) and \(aN\ge1375\). In particular the second
condition covers \(a\ge5\) at these terminals and every \(a\ge1\)
when \(N\ge3150\).
\end{lemma}

\begin{proof}
Theorem~1(a) of Laishram and Shorey~\cite{laishram-shorey} gives a prime
divisor greater than \(2H\) of the first \(H\) consecutive integers
beginning at \(aN+1\), under the stated starting-point inequality.
Lemma~\ref{lem:prime-support} then gives \(2H<C_N\).
If \(a\ge2d\), the inequality follows from \(h\le N-1\) and \(N\ge25\).

It remains to establish the starting-point inequality for
\(N\ge275\), \(aN\ge1375\), and \(a<2d\).
Put \(X=aN\), \(w=X/40\), and consider
three intervals of length \(w\) whose left endpoints are
\[
 X,\qquad X+2w+4d,\qquad X+6w+12d.
\]
Take each interval open at its left endpoint. Their union is contained
in \((X,X+E]\) with \(E=7w+12d< X/3\), since \(d<3a\).
Each interval contains a prime. Indeed \(X\ge1375\). For
\(X\le7160\), \(7160<X\le11750\), and \(11750<X<89693\), respectively,
the upper endpoint \(4X/3\) lies below \(9551,15683,155921\),
where the record prime gaps first exceed \(34,36,72\),
respectively~\cite[Table~8]{oliveira-gaps}. In each range \(w\)
exceeds the corresponding gap. For \(X\ge89693\),
Dusart~\cite[Proposition~5.4]{dusart-interval} gives a prime within
\(x/\log^3x<w\) of each left endpoint \(x\in[X,4X/3)\).

Discard any interval meeting
\[
 I_j=[(M-b)/j,(M+2d)/j]\qquad(1\le j\le4).
\]
The envelope misses \(I_1\), since \(M-b>(3/2)(1-1/N)X>4X/3\).
If \(I_j\), \(2\le j\le4\), meets the envelope, then \(M+2d>jX\).
For \(j'<j\) this implies
\[
 (M-b)/j' > X+(X-a-3d)/3>X+E,
\]
where \(E/X<7/40+36/275 < (1-10/275)/3\).
Thus at most one \(I_j\) meets the envelope. Its width is
at most \((a+3d)/2<5d/2\), smaller than the separation of the
three intervals. We obtain two primes \(p_1<p_2\) in distinct
retained intervals.

For either retained prime \(p\), we have \(p>X\ge m\) and
\(p^2>M\) because \(b<4a\). Therefore
\(y_m(k)=\lfloor(am+dk)/p\rfloor\). The avoidance of the \(I_j\)
gives the same \(q=\lfloor bm/p\rfloor\in\{1,2,3\}\) at both orders
and \((q+1)p-bm>2d\). The right endpoints of the constant-ordinate
intervals before the jumps are
\[
 k_j(m)=\left\lfloor\frac{jp-am-1}{d}\right\rfloor
 \qquad(1\le j\le q).
\]
Every positive edge before the last of these endpoints joins two such
endpoints and has slope at most \(1/\lfloor p/d\rfloor\).
This follows by subtracting their floor expressions, since an edge
of height \(J\) then has width at least \(J\lfloor p/d\rfloor\).
The edge from \(k_q(m)\) to \(m\) has height one and width
\[
 w_m=\left\lceil(bm-qp+1)/d\right\rceil
 \le \lceil p/d\rceil-2 <\lfloor p/d\rfloor.
\]
All remaining points lie on or above this final segment. Its slope
is larger than all earlier slopes, and \(w_N>w_{N-1}\) since \(b>d\).
Neither final slope is shared.

For an interior positive edge of height \(J_e\) in \(P_N\), its
width is \(w_e=J_e p/d+\epsilon_e\), with
\(|\epsilon_e|<1\). If the factor uses \(r_e\) of its roots,
the quantity \(r_eJ_e/w_e\) is an integer by
Lemma~\ref{lem:newton-alloc}. Summing over these edges yields
\[
 h=h_0+Jp/d+\epsilon,\qquad
 0\le h_0<(p-X)/d,\qquad J\in\{0,1,2\}.
\]
Here \(h_0\) counts slope-zero roots and \(J\) is at most the total
interior height \(q-1\). The error satisfies \(|\epsilon|<2\) when
\(J>0\), and is zero when \(J=0\).
For \(p_1,p_2\), the intervals of possible \(h\) with equal positive
\(J\) are disjoint because \(p_2-p_1>p_1-X+4d\), by the placement of
the intervals. Intervals for unequal \(J\) are disjoint as well,
since \(X>2E+4d\). Thus \(J=0\) at both primes. The smaller prime
comes from the first or second interval, giving
\[
 H<p_1-X\le 3w+4d<.12X .
\]
As \(X\ge1375\), this implies the required starting-point inequality.
\end{proof}

\subsection{Bounds from records of prime gaps}

For an even record gap \(G\), write \(X_G\) for the lower prime
endpoint of the first larger record gap in
\cite[Table~8]{oliveira-gaps}. Every real interval \((x,y]\)
with \(x\ge2\), \(y<X_G\) and \(y-x\ge G\) contains a prime.
Otherwise the consecutive primes straddling \(x\) have gap
greater than \(G\) and lower endpoint below \(X_G\).

\begin{lemma}[Uniform gap and degree bounds]
\label{lem:gap-caps}
Let \(G\ge52\) be an even record gap such that
\begin{equation}
 G\ge 4\sigma_N,\qquad
 G\left(\frac{N(5N-11)}6+2\right)<X_G .
 \label{eq:gap-conditions}
\end{equation}
Then \(d<G/2\). If Lemma~\ref{lem:degree-cap} applies, then also
\(H<G\).
\end{lemma}
\begin{proof}
Put \(g=G/2\). If \(d\ge g\), take \(r=\lfloor d/g\rfloor\).
Then \(g\le d/r<2g\), so \((aN/r,(aN+2d)/r]\) has length
at least \(G\). Its upper endpoint is below \(X_G\) by
\(6a<(5N-11)d\) and \eqref{eq:gap-conditions}. Its left endpoint
exceeds \(dN\), since \(s<\sigma_N\) and \(d\ge g\ge2\sigma_N\)
give \(a>d^2/\sigma_N-d\ge d^2/g\ge rd\).
A prime \(p\) in the interval divides \(rp\) among the first
\(2d\) terms of the upper factorial block, contrary to
Lemma~\ref{lem:prime-support}.

Suppose now that \(2H<C_N\) and \(H\ge G\).
Take \(r=\lfloor H/G\rfloor\). The interval
\((aN/r,(aN+H)/r]\) has length at least \(G\), and its upper
endpoint is less than
\[
 2G\left(\frac{N(5N-11)}{6h}+1\right)
 \le G\left(\frac{N(5N-11)}6+2\right)<X_G .
\]
The left endpoint exceeds two since \(r\le H/G<dN/G\) and
\(aG/d>G/3\).
If \(C_N=\sqrt M\), then \(aN/r\ge aGN/H>2aG\sqrt M/b>\sqrt M\).
If \(C_N=a+2d-1\), then \(2H<C_N\) and \(h\ge2\) give \(a>2d\).
Since \(d<G/2\), we have \(r<N/2\) and \(aN/r>2a>C_N\).
Thus in both cases the interval supplies a prime \(p>C_N\)
dividing one of the first \(H\) terms. This contradicts
Lemma~\ref{lem:prime-support}.
\end{proof}

Table~\ref{tab:gap-bounds} gives choices of \(G\) for
Lemma~\ref{lem:gap-caps}. Use the first applicable row. The
inequalities in \eqref{eq:gap-conditions} follow by checking
the right endpoint of each row, since \(2\log N+3\) is an
increasing upper bound for \(\sigma_N\).
\begin{table}[ht]
\centering
\begin{tabular}{r r r}
\hline
Upper terminal \(N\) & \(G\) & \(X_G\)\\
\hline
50 & 72 & \(155\,921\)\\
100 & 114 & \(1\,349\,533\)\\
300 & 154 & \(17\,051\,707\)\\
\(1\,000\) & 222 & \(189\,695\,659\)\\
\(2\,000\) & 282 & \(1\,294\,268\,491\)\\
\(6\,000\) & 354 & \(10\,726\,904\,659\)\\
\(10\,000\) & 456 & \(42\,652\,618\,343\)\\
\(19\,999\) & 468 & \(182\,226\,896\,239\)\\
\(32\,767\) & 516 & \(461\,690\,510\,011\)\\
\(65\,535\) & 602 & \(2\,614\,941\,710\,599\)\\
\(3\,269\,017\) & 1132 & \(43\,841\,547\,845\,541\,059\)\\
\hline
\end{tabular}
\caption{Prime-gap bounds used to control \(d\) and \(dh\). The record
endpoints are from \cite[Table~8]{oliveira-gaps}.}
\label{tab:gap-bounds}
\end{table}

\begin{lemma}[Refinement using the actual factorial position]
\label{lem:local-gap}
Assume \(2H<C_N\) and \(H\le H_0\). Let \(G\ge52\) be an even record
gap. If \(H_0/G<N/2\) and \(aN+H_0<X_G\), then
\(H<G\) and \(d<G/2\).
\end{lemma}
\begin{proof}
If \(H\ge G\), the interval with \(r=\lfloor H/G\rfloor\) in the
second part of Lemma~\ref{lem:gap-caps} has upper endpoint below
\(X_G\) by assumption, as well as length at least \(G\). The two
comparisons with \(C_N\) in that proof still hold, using
\(r\le H_0/G<N/2\) in the second case. The resulting prime gives
the same contradiction.
\end{proof}

In particular \(H\le281\) and \(d\le140\) for
\(1001\le N\le3295\) whenever Lemma~\ref{lem:degree-cap} applies.
Starting with \(G=354\) from
Table~\ref{tab:gap-bounds}, apply Lemma~\ref{lem:local-gap}
successively with
\[
\begin{array}{r|r|r}
G&\text{upper terminal }N&X_G\\ \hline
320&5000&3\,842\,610\,773\\
292&4150&2\,300\,942\,549\\
288&3467&1\,453\,168\,141\\
282&3295&1\,294\,268\,491
\end{array}
\]
At each step \(aN+H_0<(5/6)D_0N^2+H_0<X_G\), where \(d\le D_0\)
and \(H\le H_0\) are the preceding bounds.

\subsection{Exact computation in the finite range}

\begin{proposition}
\label{prop:finite}
Under the support reduction of Lemma~\ref{lem:support-reduction},
the consecutive polynomials are coprime for \(a\ge5\), \(d\ge2\),
and \(25\le N<3150\).
\end{proposition}
\begin{proof}
Use \(d^2<19b\), which follows from \(\sigma_N<19\) in this
range, and use the gap caps
\[
\begin{array}{c|ccccc}
N&25\text{--}50&51\text{--}100&101\text{--}300&301\text{--}1000&
1001\text{--}3149\\
d\le&35&56&76&110&140 .
\end{array}
\]
These come from Lemmas~\ref{lem:gap-caps} and \ref{lem:local-gap}.
For \(2200\le N<3150\), the analytic exclusions in
Appendix~\ref{app:middle-tail} also require \(b>N\) and
\(97d<3a\).

For each \(N,a\), the inequalities \(6a<(5N-11)d\), \(2\le d<3a\),
the displayed cap, and the two additional inequalities when
\(N\ge2200\) give an integer interval \([d_0,d_1]\) for \(d\).
When \(N\ge275\) or \(a\ge2d\), Lemma~\ref{lem:degree-cap}
reduces the bound in Lemma~\ref{lem:prime-support} to
\(\min(dN,C_N)\), independent of \(h\). Otherwise use \(dN\).
The resulting bound is nondecreasing in \(d\), including where
it switches to \(dN\). So a prime exceeding the bound at \(d_1\)
in either first \(2d_0\)-term block eliminates the whole interval.
The remaining gaps are tested individually with their first
\(2d\) terms and the conditions \(a<d(N-2)\), \(d^2<19b\).

For each remaining support, try the possible even degrees
\(2\le h<N-1\). Test the starting-point inequality directly
when \(N<275\). Impose the degree bound when it applies, and
test the full first \(H\) terms with
Lemma~\ref{lem:prime-support} in both orders.

The remaining degrees are tested at odd primes \(aN<p\le M\).
Here \(p>m,d\), \(p^2>M\), and \(am<p\). The ordinate is
\(\lfloor(am+dk)/p\rfloor\) and \(q_m=\lfloor bm/p\rfloor\le3\).
Thus the lower hull uses at most five points
\[
 (0,0),\qquad
 \left(\left\lfloor\frac{jp-am-1}{d}\right\rfloor,j-1\right)
 \quad(1\le j\le q_m),\qquad (m,q_m).
\]
Every omitted point is to the left of a retained point of the
same ordinate and lies on or above the retained lower hull,
which is nondecreasing. Compute the allowed degree sum from
Lemma~\ref{lem:newton-alloc} at each prime and intersect it with
the surviving degrees. For example \(N=25,a=5,d=8,p=127\)
gives hull slopes \(1/16,1/8\) and \(1/16,1/9\) at orders
24 and 25, respectively. Only degrees zero and sixteen are
allowed at this prime.

The supplementary verifier exhausts these integer intervals.
The largest factorial argument is \(1\,156\,520\,634\).
A segmented sieve stores smallest odd prime factors through that limit.
Trial divisions from this table and cross multiplication of slopes
perform the tests above. In prime caps the square-root threshold
is replaced by its integer part, checked by squaring. The initial prime-cap test leaves
\(21\,583\) groups \((N,a)\). The individual gap test leaves
\(74\,362\) supports. The degree and prime-block tests exclude
\(62\,593\), and the Newton allocations exclude \(11\,705\).
For the remaining \(64\) supports, the verifier forms both
monic reciprocals modulo the prime \(1\,000\,000\,007\) and
checks by Euclidean division that their gcd is one. The
modulus is checked by trial division. A monic integral common
divisor would retain its degree on reduction, so these
computations prove rational coprimality.
\end{proof}
